% ============================================================
%  Chapter 4 -- Connectedness, Discontinuities, Monotonic
%               Functions, Infinite Limits: 4.22-4.34
% ============================================================
\rsection{Continuity and Connectedness}

\begin{signpost}[One theorem, one corollary, one warning]
The connectedness investment of 2.45--2.47 pays out in a single exchange:
continuous images of connected sets are connected (4.22), and since the
connected subsets of $\R$ are exactly the intervals (2.47), the intermediate
value theorem falls out in three citations (4.23). The warning is 4.24: the
intermediate value \emph{property} does not characterise continuity ---
$\sin(1/x)$ has it while being badly discontinuous. The converse a beginner
expects is false.
\end{signpost}

\ritem{4.22}{Theorem}
\emph{If $f$ is a continuous mapping of a metric space $X$ into a metric
space $Y$, and if $E$ is a connected subset of $X$, then $f(E)$ is
connected.}

\begin{proof}
Assume, on the contrary, that $f(E) = A \cup B$, where $A$ and $B$ are
nonempty separated subsets of $Y$. Put $G = E \cap f^{-1}(A)$,
$H = E \cap f^{-1}(B)$.

Then $E = G \cup H$, and neither $G$ nor $H$ is empty.

Since $A \subset \overline{A}$ (the closure of $A$), we have
$G \subset f^{-1}(\overline{A})$; the latter set is closed, since $f$ is
continuous; hence $\overline{G} \subset f^{-1}(\overline{A})$. It follows
that $f(\overline{G}) \subset \overline{A}$. Since $f(H) = B$ and
$\overline{A} \cap B$ is empty, we conclude that $\overline{G} \cap H$ is
empty.

The same argument shows that $G \cap \overline{H}$ is empty. Thus $G$ and $H$
are separated. This is impossible if $E$ is connected.\kn{The same transport
strategy as 4.14, run contrapositively: a separation of the image pulls back
to a separation of the domain. The only technical point is that
``separated'' involves \emph{closures}, and this is where the corollary to
4.8 earns its place --- $f^{-1}(\overline A)$ is closed, so it swallows
$\overline G$, keeping the pulled-back pieces apart. Images transport
connectedness and compactness; inverse images transport openness and
closedness. Keep the direction ledger straight.}
\end{proof}

\ritem{4.23}{Theorem}
\emph{Let $f$ be a continuous real function on the interval $[a,b]$. If
$f(a) < f(b)$ and if $c$ is a number such that $f(a) < c < f(b)$, then there
exists a point $x \in (a,b)$ such that $f(x) = c$.}

A similar result holds, of course, if $f(a) > f(b)$. Roughly speaking, the
theorem says that a continuous real function assumes all intermediate values
on an interval.

\begin{proof}
By Theorem 2.47, $[a,b]$ is connected; hence Theorem 4.22 shows that
$f([a,b])$ is a connected subset of $\R^1$, and the assertion follows if we
appeal once more to Theorem 2.47.\kn{Three citations and no computation:
2.47 converts ``interval'' to ``connected,'' 4.22 pushes connectedness
through $f$, and 2.47 converts back to ``interval'' --- which by definition
contains everything between $f(a)$ and $f(b)$. Chapter 1's Theorem 1.21 was
this argument done by hand, before any of the three tools existed; compare
the box there. Exercise 4.14 (fixed points) is the standard application.}
\end{proof}

\begin{center}
\begin{tikzpicture}[x=1mm, y=1mm]
  \draw[axis] (-2,0) -- (78,0);
  \draw[axis] (0,-2) -- (0,32);
  % the curve
  \draw[thin] plot[smooth, tension=0.9]
    coordinates {(6,5) (18,14) (30,8) (44,22) (56,18) (68,27)};
  \node[ptfill, minimum size=2.4pt] at (6,5) {};
  \node[ptfill, minimum size=2.4pt] at (68,27) {};
  \node[mlbl, anchor=north] at (6,-1.5) {$a$};
  \node[mlbl, anchor=north] at (68,-1.5) {$b$};
  \draw[guide] (6,5) -- (6,0);
  \draw[guide] (68,27) -- (68,0);
  % level c
  \draw[guide] (0,16) -- (76,16);
  \node[mlbl, anchor=east] at (-1.5,16) {$c$};
  % crossings
  \foreach \x in {21.5, 36.5, 52}{\node[ptopen, minimum size=2.8pt] at (\x,16) {};}
  \node[lbl, anchor=south] at (37,29) {every level between $f(a)$ and $f(b)$};
  \node[lbl, anchor=south] at (37,25.5) {is crossed --- possibly several times};
\end{tikzpicture}
\end{center}
\figcap{The intermediate value theorem: a connected image cannot skip the
level $c$. The theorem asserts at least one crossing and says nothing about
how many, nor where.}

\ritem{4.24}{Remark}
At first glance, it might seem that Theorem 4.23 has a converse. That is, one
might think that if for any two points $x_1 < x_2$ and for any number $c$
between $f(x_1)$ and $f(x_2)$ there is a point $x$ in $(x_1,x_2)$ such that
$f(x) = c$, then $f$ must be continuous.

That this is not so may be concluded from Example 4.27(d).

\rsection{Discontinuities}

\begin{signpost}[A taxonomy, then a zoo]
Failure of continuity comes in two grades. If both one-sided limits exist,
the failure is a bookkeeping error --- a jump, or a wrong value --- and is
called \emph{simple} (first kind). If a one-sided limit fails to exist, the
function oscillates or scatters, and the discontinuity is of the second kind.
Example 4.27 exhibits all the species; (a) and (b) are the standard
pathologies built from $\Q$, and (d) is the oscillator that killed the
converse of the intermediate value theorem in 4.24. The taxonomy exists for
the next section: monotonic functions will be shown to admit only the first
kind, and only countably many.
\end{signpost}

\noindent
If $x$ is a point in the domain of definition of the function $f$ at which
$f$ is not continuous, we say that $f$ is \emph{discontinuous at $x$}, or
that $f$ has a \emph{discontinuity} at $x$. If $f$ is defined on an interval
or on a segment, it is customary to divide discontinuities into two types.
Before giving this classification, we have to define the \emph{right-hand}
and the \emph{left-hand limits} of $f$ at $x$, which we denote by $f(x+)$ and
$f(x-)$, respectively.

\ritem{4.25}{Definition}
Let $f$ be defined on $(a,b)$. Consider any point $x$ such that
$a \le x < b$. We write
\[ f(x+) = q \]
if $f(t_n) \to q$ as $n \to \infty$, for all sequences $(t_n)$ in $(x,b)$
such that $t_n \to x$. To obtain the definition of $f(x-)$, for
$a < x \le b$, we restrict ourselves to sequences $(t_n)$ in $(a,x)$.

It is clear that for any point $x$ of $(a,b)$, $\lim_{t\to x} f(t)$ exists if
and only if
\[ f(x+) = f(x-) = \lim_{t\to x} f(t). \]

\ritem{4.26}{Definition}
Let $f$ be defined on $(a,b)$. If $f$ is discontinuous at a point $x$, and if
$f(x+)$ and $f(x-)$ exist, then $f$ is said to have a \emph{discontinuity of
the first kind}, or a \emph{simple discontinuity}, at $x$. Otherwise the
discontinuity is said to be \emph{of the second kind}.

There are two ways in which a function can have a simple discontinuity:
either $f(x+) \ne f(x-)$ [in which case the value $f(x)$ is immaterial], or
$f(x+) = f(x-) \ne f(x)$.

\ritem{4.27}{Examples}
\begin{enumerate}[label=(\alph*), leftmargin=2.2em, itemsep=4pt, topsep=4pt]
  \item Define
        \[ f(x) = \begin{cases} 1 & (x \text{ rational}),\\
           0 & (x \text{ irrational}). \end{cases} \]
        Then $f$ has a discontinuity of the second kind at every point $x$,
        since neither $f(x+)$ nor $f(x-)$ exists.\kn{The Dirichlet function,
        and the standard witness that ``discontinuous everywhere'' is
        possible. Both $\Q$ and its complement are dense (1.20(b) and
        Exercise 1.1), so every one-sided approach can be made along either
        set, producing both $0$ and $1$ as subsequential limits --- the
        refutation direction of the bridge 4.2. It returns in Chapter 6 as
        the function that is not Riemann-integrable (Exercise 6.4).}
  \item Define
        \[ f(x) = \begin{cases} x & (x \text{ rational}),\\
           0 & (x \text{ irrational}). \end{cases} \]
        Then $f$ is continuous at $x = 0$ and has a discontinuity of the
        second kind at every other point.\kn{Continuity at exactly one point
        --- worth a pause, since beginners often believe continuity
        somewhere forces decency nearby. The squeeze $|f(x)| \le |x|$ gives
        continuity at $0$ no matter how wild the values elsewhere.}
  \item Define
        \[ f(x) = \begin{cases} x+2 & (-3 < x < -2),\\
           -x-2 & (-2 \le x < 0),\\ x+2 & (0 \le x < 1). \end{cases} \]
        Then $f$ has a simple discontinuity at $x = 0$ and is continuous at
        every other point of $(-3,1)$.
  \item Define
        \[ f(x) = \begin{cases} \sin\dfrac1x & (x \ne 0),\\[4pt]
           0 & (x = 0). \end{cases} \]
        Since neither $f(0+)$ nor $f(0-)$ exists, $f$ has a discontinuity of
        the second kind at $x = 0$. We have not yet shown that $\sin x$ is a
        continuous function. If we assume this result for the moment,
        Theorem 4.7 implies that $f$ is continuous at every point
        $x \ne 0$.\kd{The oscillator that closes 4.24: on any interval
        $(0,\delta)$ it sweeps all of $[-1,1]$, so it has the full
        intermediate value property, yet it is discontinuous at $0$. The
        intermediate value property constrains what a function \emph{skips},
        not how it \emph{settles} --- and settling is what continuity is.
        Exercise 4.14's fixed-point result survives for such functions;
        their integrability will be delicate in Chapter 6.}
\end{enumerate}

\begin{center}
\begin{tikzpicture}[x=1mm, y=1mm]
  % ---- simple discontinuity (jump), from 4.27(c)
  \draw[axis] (-2,0) -- (40,0);
  \draw[axis] (0,-2) -- (0,24);
  \draw[thin] (2,4) -- (16,11);
  \node[ptopen] at (16,11) {};
  \node[ptfill, minimum size=2.6pt] at (16,19) {};
  \draw[thin] (16,19) -- (34,23);
  \draw[guide] (16,0) -- (16,19);
  \node[mlbl, anchor=north] at (16,-1.5) {$x$};
  \node[mlbl, anchor=east] at (14.5,11) {$f(x-)$};
  \node[mlbl, anchor=south east] at (15,19.5) {$f(x+)$};
  \node[lbl, anchor=north] at (18,-6) {first kind: both one-sided};
  \node[lbl, anchor=north] at (18,-9.5) {limits exist, disagree};
  % ---- second kind: oscillation sin(1/x)
  \begin{scope}[xshift=54mm]
  \draw[axis] (-2,0) -- (42,0);
  \draw[axis] (0,-2) -- (0,24);
  \draw[guide] (0,21) -- (40,21);
  \draw[guide] (0,3) -- (40,3);
  \draw[thin] plot[smooth, tension=0.8] coordinates
    {(3,12) (4,20.4) (5,3.6) (6.2,20.4) (7.6,3.6) (9.4,20.4) (11.6,3.6)
     (14.4,20.4) (18,3.6) (23,20.4) (30,6) (40,9.5)};
  \node[mlbl, anchor=north] at (3,-1.5) {$x$};
  \draw[guide] (3,0) -- (3,3);
  \node[lbl, anchor=north] at (20,-6) {second kind: oscillation};
  \node[lbl, anchor=north] at (20,-9.5) {denser and denser, no limit};
  \end{scope}
\end{tikzpicture}
\end{center}
\figcap{The two kinds. Left, a jump (4.27(c)): the function arrives at
different platforms from the two sides. Right, $\sin(1/x)$ near $0$
(4.27(d)): the one-sided limits do not exist at all, because every
neighbourhood of $x$ sweeps the whole band.}

\rsection{Monotonic Functions}

\noindent
We shall now study those functions which never decrease (or never increase)
on a given segment.

\ritem{4.28}{Definition}
Let $f$ be real on $(a,b)$. Then $f$ is said to be \emph{monotonically
increasing} on $(a,b)$ if $a < x < y < b$ implies $f(x) \le f(y)$. If the
last inequality is reversed, we obtain the definition of a
\emph{monotonically decreasing} function. The class of \emph{monotonic}
functions consists of both the increasing and the decreasing functions.

\ritem{4.29}{Theorem}
\emph{Let $f$ be monotonically increasing on $(a,b)$. Then $f(x+)$ and
$f(x-)$ exist at every point $x$ of $(a,b)$. More precisely,}
\begin{equation}
  \sup_{a<t<x} f(t) = f(x-) \le f(x) \le f(x+) = \inf_{x<t<b} f(t).
  \tag{4.25}
\end{equation}
\emph{Furthermore, if $a < x < y < b$, then}
\begin{equation}
  f(x+) \le f(y-). \tag{4.26}
\end{equation}
\emph{Analogous results evidently hold for monotonically decreasing
functions.}

\begin{proof}
By hypothesis, the set of numbers $f(t)$, where $a < t < x$, is bounded above
by the number $f(x)$, and therefore has a least upper bound which we shall
denote by $A$. Evidently $A \le f(x)$. We have to show that $A = f(x-)$.

Let $\varepsilon > 0$ be given. It follows from the definition of $A$ as a
least upper bound that there exists $\delta > 0$ such that $a < x-\delta < x$
and
\begin{equation}
  A - \varepsilon < f(x-\delta) \le A. \tag{4.27}
\end{equation}
Since $f$ is monotonic, we have
\begin{equation}
  f(x-\delta) \le f(t) \le A \qquad (x-\delta < t < x). \tag{4.28}
\end{equation}
Combining (4.27) and (4.28), we see that
\[ |f(t)-A| < \varepsilon \qquad (x-\delta < t < x). \]
Hence $f(x-) = A$.\kn{The bounded-implies-sup pattern of the strategy index
(pattern 3), in its function-limit costume: monotonicity makes the left
values bounded, the supremum $A$ exists by 1.19, and clause (ii) of 1.8
plants a term within $\varepsilon$; monotonicity then drags the whole
approach into the gap --- exactly the mechanism of 3.14, with $t \uparrow x$
in place of $n \to \infty$.}

The second half of (4.25) is proved in precisely the same way.

Next, if $a < x < y < b$, we see from (4.25) that
\begin{equation}
  f(x+) = \inf_{x<t<b} f(t) = \inf_{x<t<y} f(t). \tag{4.29}
\end{equation}
The last equality is obtained by applying (4.25) to $(a,y)$ in place of
$(a,b)$. Similarly,
\begin{equation}
  f(y-) = \sup_{a<t<y} f(t) = \sup_{x<t<y} f(t). \tag{4.30}
\end{equation}
Comparison of (4.29) and (4.30) gives (4.26).
\end{proof}

\ritem{}{Corollary}
\emph{Monotonic functions have no discontinuities of the second kind.}

This corollary implies that every monotonic function is discontinuous at a
countable set of points at most. Instead of appealing to the general theorem
whose proof is sketched in Exercise 4.17, we give here a simple proof which
is applicable to monotonic functions.

\ritem{4.30}{Theorem}
\emph{Let $f$ be monotonic on $(a,b)$. Then the set of points of $(a,b)$ at
which $f$ is discontinuous is at most countable.}

\begin{proof}
Suppose, for the sake of definiteness, that $f$ is increasing, and let $E$ be
the set of points at which $f$ is discontinuous.

With every point $x$ of $E$ we associate a rational number $r(x)$ such that
\[ f(x-) < r(x) < f(x+). \]
Since $x_1 < x_2$ implies $f(x_1+) \le f(x_2-)$, we see that
$r(x_1) \ne r(x_2)$ if $x_1 \ne x_2$.

We have thus established a 1-1 correspondence between the set $E$ and a
subset of the set of rational numbers. The latter, as we know, is
countable.\kn{A counting proof with real geometry in it: by 4.29 every
discontinuity of an increasing $f$ is a \emph{jump}, the jump intervals
$(f(x-),f(x+))$ are pairwise disjoint by (4.26), and disjoint intervals can
each be tagged with their own rational --- density of $\Q$ (1.20(b)) supplies
the tag, countability of $\Q$ (2.13) does the counting. The same
disjoint-intervals-are-countable argument proved Exercise 2.29. (Choosing
$r(x)$ for each $x$ is another silent appeal to choice; here it can be
avoided --- take the rational with least index in a fixed enumeration.)}
\end{proof}

\ritem{4.31}{Remark}
It should be noted that the discontinuities of a monotonic function need not
be isolated. In fact, given any countable subset $E$ of $(a,b)$, which may
even be dense, we can construct a function $f$, monotonic on $(a,b)$,
discontinuous at every point of $E$, and at no other point of $(a,b)$.

To show this, let the points of $E$ be arranged in a sequence $(x_n)$,
$n = 1,2,3,\dots$. Let $(c_n)$ be a sequence of positive numbers such that
$\Sigma c_n$ converges. Define
\begin{equation}
  f(x) = \sum_{x_n<x} c_n \qquad (a < x < b). \tag{4.31}
\end{equation}
The summation is to be understood as follows: Sum over those indices $n$ for
which $x_n < x$. If there are no points $x_n$ to the left of $x$, the sum is
empty; following the usual convention, we define it to be zero. Since (4.31)
converges absolutely, the order in which the terms are arranged is
immaterial.\kd{Chapter 3's rearrangement theorem (3.55) quietly holding up
the definition: the indices with $x_n < x$ arrive in no particular order, so
the sum is well defined \emph{only because} the convergence is absolute. A
conditionally convergent $(c_n)$ would make (4.31) meaningless. This jump
function is the prototype for the integrators of Chapter 6 (Rudin points to
6.16), where its jumps become point masses.}

We leave the verification of the following properties of $f$ to the reader:
\begin{enumerate}[label=(\alph*), leftmargin=2.2em, itemsep=1pt, topsep=3pt]
  \item $f$ is monotonically increasing on $(a,b)$;
  \item $f$ is discontinuous at every point of $E$; in fact,
        \[ f(x_n+) - f(x_n-) = c_n. \]
  \item $f$ is continuous at every other point of $(a,b)$.
\end{enumerate}
Moreover, it is not hard to see that $f(x-) = f(x)$ at all points of $(a,b)$.
If a function satisfies this condition, we say that $f$ is \emph{continuous
from the left}. If the summation in (4.31) were taken over all indices $n$
for which $x_n \le x$, we would have $f(x+) = f(x)$ at every point of
$(a,b)$; that is, $f$ would be \emph{continuous from the right}.

Functions of this sort can also be defined by another method; for an example
we refer to Theorem 6.16.

\rsection{Infinite Limits and Limits at Infinity}

\begin{signpost}[The extended reals, cashed in]
Two definitions and a theorem, closing the chapter by making
$\lim_{x\to\infty}$ and ``$f(t) \to +\infty$'' legitimate. The device is the
one hinted at in 1.23: give $\pm\infty$ \emph{neighbourhoods}, and Definition
4.1 restates itself with neighbourhoods in place of distances --- no new
theory needed, and 4.34 lists which limit laws survive, with precisely the
exceptions catalogued as undefined at 1.23.
\end{signpost}

\noindent
To enable us to operate in the extended real number system, we shall now
enlarge the scope of Definition 4.1 by reformulating it in terms of
neighborhoods.

For any real number $x$, we have already defined a neighborhood of $x$ to be
any segment $(x-\delta,\,x+\delta)$.

\ritem{4.32}{Definition}
For any real $c$, the set of real numbers $x$ such that $x > c$ is called a
\emph{neighborhood of $+\infty$} and is written $(c,+\infty)$. Similarly, the
set $(-\infty,c)$ is a neighborhood of $-\infty$.

\ritem{4.33}{Definition}
Let $f$ be a real function defined on $E \subset \R$. We say that
\[ f(t) \to A \text{ as } t \to x, \]
where $A$ and $x$ are in the extended real number system, if for every
neighborhood $U$ of $A$ there is a neighborhood $V$ of $x$ such that
$V \cap E$ is not empty, and such that $f(t) \in U$ for all
$t \in V \cap E$, $t \ne x$.\kn{One definition now covers nine cases ($x$
and $A$ each finite, $+\infty$, or $-\infty$), because ``neighbourhood''
absorbs the difference: an $\varepsilon$-ball around a real number, a ray
beyond $c$ for $\pm\infty$. This is the topologist's trick --- rephrase
metrically-stated definitions in terms of neighbourhoods and they generalise
by themselves. Theorem 3.15's sequence version was the same move.}

A moment's consideration will show that this coincides with Definition 4.1
when $A$ and $x$ are real.

The analogue of Theorem 4.4 is still true, and the proof offers nothing new.
We state it, for the sake of completeness.

\ritem{4.34}{Theorem}
\emph{Let $f$ and $g$ be defined on $E \subset \R$. Suppose}
\[ f(t) \to A, \qquad g(t) \to B \qquad\text{as } t \to x. \]
\emph{Then}
\begin{enumerate}[label=(\alph*), leftmargin=2.2em, itemsep=1pt, topsep=3pt]
  \item $f(t) \to A'$ \emph{implies} $A' = A$,
  \item $(f+g)(t) \to A+B$,
  \item $(fg)(t) \to AB$,
  \item $(f/g)(t) \to A/B$,
\end{enumerate}
\emph{provided the right members of (b), (c), and (d) are defined.}

Note that $\infty-\infty$, $0\cdot\infty$, $\infty/\infty$, and $A/0$ are not
defined (see Definition 1.23).

\begin{recap}[Chapter 4 in one paragraph]
Continuity is one property in three costumes: $\varepsilon$--$\delta$ at a
point (4.5), preservation of sequential limits (4.2 with 4.6), and open
inverse images (4.8) --- the last being the form that proves theorems. What
continuity transports is the content of the chapter: compactness pushes
forward (4.14), and with it boundedness, attained extrema, continuous
inverses, and uniform continuity (4.15--4.19); connectedness pushes forward
(4.22), and with it the intermediate value theorem (4.23). The
counterexamples 4.20, 4.21, 4.24 and 4.27 mark the exact boundaries: nothing
survives the loss of compactness except by accident, and the intermediate
value property is weaker than continuity. Monotonic functions are the tame
class --- one-sided limits always exist (4.29), so only jumps occur, and only
countably many (4.30), though possibly on a dense set (4.31). Chapter 6 will
integrate against exactly such jump functions, and Chapter 5's derivative
inherits every tool built here.
\end{recap}
